%%%%%%%%%%%%%%%%%%%%%%%%%%%%%%%%%%%%%%%%%
% Protocolo Laboratorio de Fisicoqu\'{i}mica Te\'{o}rica 
% LaTeX Template
% Version 0.1 (10/12/17)
%
% Part of this template was downloaded from:
% http://www.LaTeXTemplates.com
%
% Edgar Garc\'{i}a Juarez (b613) edgargj.16@gmail.com
%%%%%%%%%%%%%%%%%%%%%%%%%%%%%%%%%%%%%%%%%
%-----------------------------------------------------------------------------% 
% Packages
%-----------------------------------------------------------------------------%
\documentclass[11pt]{article} 	
\usepackage[utf8]{inputenc}
\usepackage[spanish]{babel}
\usepackage{latexsym,amssymb,amsfonts}
\usepackage[usenames]{color}
\usepackage{float}
\usepackage{graphicx}
\usepackage{verbatim}

\usepackage{amsmath,mathtools}
\usepackage[font={small},format=plain,labelfont=bf]{caption} % Size and format of foot of figures and tables
\usepackage{color}
\usepackage{hyperref}
%-----------------------------------------------------------------------------%
% Margin settings 
%-----------------------------------------------------------------------------%
\usepackage{geometry}
\geometry{
	paper=letterpaper, % Paper type
	inner=2.5cm, % Inner margin
	outer=2.5cm, % Outer margin
	top=3.5cm, % Top margin
	bottom=3.5cm, % Bottom margin
}
%-----------------------------------------------------------------------------% 
% New Commands
%-----------------------------------------------------------------------------%
\renewcommand{\tablename}{Tabla} % Name of table foot 
\renewcommand{\baselinestretch}{1.5} % space between lines 1.5
\newcommand\tab[1][0.5cm]{\hspace*{#1}}
\newcommand\vtab[1][0.5cm]{\vspace*{#1}}
%-----------------------------------------------------------------------------% 
%-----------------------------------------------------------------------------%
% Document
%-----------------------------------------------------------------------------%
%-----------------------------------------------------------------------------% 

\begin{document}
%-----------------------------------------------------------------------------%
% Cover Page
%-----------------------------------------------------------------------------% 
\pagestyle{empty} 
\phantom{a}
\vspace{-1.4cm}
\begin{center}
\includegraphics[width=4cm]{buap.pdf}
%\includegraphics[width=4cm]{fcq.pdf}

\Large{\textbf{Benem\'{e}rita Universidad Aut\'{o}noma de Puebla}\\
\rule{150mm}{0.1mm}\\
Facultad de Ciencias Qu\'{i}micas\\
Centro de Qu\'{i}mica - Instituto de Ciencias\\
Posgrado en Ciencias Qu\'{i}micas}\\
\rule{150mm}{0.1mm}

\vtab[.1cm]
{\scshape\large \textbf{La dimensi\'{o}n fractal de masa como medida caracterizadora\\de la estructura de las prote\'{i}nas}}\\
\vtab[.1cm]
\Large{Protocolo de tesis de Maestr\'{i}a}\\
\vtab[0.2cm]

Lic. en Qu\'{i}mica \\
Edgar Garc\'{i}a Ju\'{a}rez \\
\vtab[0.5cm]
\begin{tabular}{cc}
\vtab[10mm]
Director de tesis & Co-directora de tesis \\
\rule{50mm}{0.1mm} & \rule{50mm}{0.1mm} \\
Dr. J. M. Solano Altamirano & Dra. Viridiana Vargas Castro \\ 
FCQ - BUAP & FCQ - BUAP \\ 
\end{tabular} 
\vtab[1cm] \\
{\today}
\end{center}

\clearpage

\section{Introducci\'{o}n}

A pesar de su aparente simplicidad, el problema de cuantificar las diferencias entre dos estruc-
turas de la misma prote\'{i}na no es trivial y continúa evolucionando. Debido a que el aspecto y distribuci\'{o}n de las prote\'{i}nas es aparentemente desordenado, la dimensi\'{o}n fractal  puede sintetizar toda la informaci\'{o}n que contienen estos sistemas. 



\section{Antecedentes}

\subsection{Fractal y dimesi\'{o}n fractal}
\label{subsec:subseccion2.1}


Existen objetos que tienen la propiedad de que su geometr\'{i}a tenga una apariencia fragmentada o con una morfolog\'{i}a aparentemente irregular y que no cambia a cualquiera que sea la escala a la que se analice; a estos objetos se les conoce como fractales. Los fractales est\'{a}n presentes en un gran n\'{u}mero de sistemas complejos que van desde agregados microsc\'{o}picos hasta c\'{u}mulos de galaxias, por lo tanto, existen objetos que se pueden describir en t\'{e}rminos de su geometr\'{i}a fractal \cite{Vicsek1992}. 

La autosimilitud o invariancia de escala es la propiedad m\'{a}s importante en la geometr\'{i}a fractal porque puede modelar la complejidad de cualquier objeto. La dimensi\'{o}n fractal \textbf{(D)} es el \'{i}ndice num\'{e}rico que cuantifica esta complejidad, midiendo la tasa de adici\'{o}n de detalles estructurales con una mayor ampliaci\'{o}n.

En la literatura cient\'{i}fica puede encontrase una gran variedad de m\'{e}todos para determinar la dimensi\'{o}n fractal, algunos de ellos son la dimensi\'{o}n fractal Hausdorff-Besicovitch, la dimensi\'{o}n fractal Minkowski y la dimesi\'{o}n de fractal por conteo de cajas. La dimensi\'{o}n fractal puede representarse como:

\begin{equation}
 D = \frac{\ln(N)}{\ln(1/\varepsilon)}
\end{equation}

Donde \(N\) es el n\'{u}mero de unidades constitutivas en la generaci\'{o}n actual que han emergido de una unidad constitutiva en la generaci\'{o}n anterior, cuando se avanza en la construcci\'{o}n del patr\'{o}n una etapa m\'{a}s adelante, o el n\'{u}mero de iniciadores reducidos que constituyen el generador. Se puede generalizar la ecuaci\'{o}n anterior para examinar la autosimilitud de varios patrones aleatorios y determinar \(D\).
 
\subsection{Determinacion de la dimensi\'{o}n fractal}
\label{subsec:subseccion2.2}


La determinacion de la dimensi\'{o}n fractal de estructuras en crecimiento, normalmente resulta de la aplicaci\'{o}n directa de las definiciones para D, sin embargo, esto es ineficaz. Es necesario medir o calcular cantidades que se puede demostrar que est\'{a}n relacionadas con la dimensi\'{o}n fractal de los objetos. Existen  tres tipos de enfoques principales para la determinaci\'{o}n de estas cantidades: (i) el experimental, (ii) el te\'{o}rico y (iii) el inform\'{a}tico \cite{Vicsek1992}. El enfoque inform\'{a}tico se obtienen por dos m\'{e}todos principales: 


\begin{enumerate}
\item Digitalizando im\'{a}genes tomadas de objetos producidos en experimentos.
\item Usando procedimientos num\'{e}ricos utilizados para simular diversos fen\'{o}menos de crecimiento.
\end{enumerate}
 
  En el segundo caso, los datos generados num\'{e}ricamente suelen producirse mediante variaciones del m\'{e}todo de Monte Carlo. A continuaci\'{o}n, discutiremos c\'{o}mo medir D para un solo objeto. Para hacer las estimaciones m\'{a}s precisas, generalmente se calcula la dimensi\'{o}n fractal para muchos grupos y se promedia sobre los resultados. Quiz\'{a}s, el m\'{e}todo m\'{a}s simple es usar la definici\'{o}n de  D como se indica en \ref{subsec:subseccion2.1}. Por lo tanto, la dimensi\'{o}n fractal se puede obtener ajustando una l\'{i}nea recta a la parte asint\'{o}tica de los datos de  N(R), por ejemplo, usando el m\'{e}todo de m\'{i}nimos cuadrados \cite{Vicsek1992}.

\subsection{M\'{e}todo de conteo de cajas}

Este m\'{e}todo consiste en dividir el espacio que contiene un patr\'{o}n dado en celdas o p\'{i}xeles de longitud (E), despu\'{e}s se cuenta el n\'{u}mero de celdas o p\'{i}xeles que existen \(N(E)\). El n\'{u}mero \(N(E)\) se mide variando el valor de \(E\) y traz\'{a}ndolo en escala logar\'{i}tmica frente a \(E\), esta es una forma pr\'{a}ctica de comprobar la autosimilitud de un patr\'{o}n determinado. Si los puntos se ajustan a una l\'{i}nea recta, se puede concluir que el patr\'{o}n dado es autosimilar, y la dimensi\'{o}n fractal \(D\) se obtiene del valor de la pendiente de la l\'{i}nea \cite{Mustafa1996}. 

\subsection{Relaci\'{o}n entre Masa y radio}

La relaci\'{o}n masa-radio es \'{u}til para estimar la dimensi\'{o}n de objetos similares a grupos (redes, vasos sangu\'{i}neos, grupos de agregaci\'{o}n). Consiste en seleccionar un punto de origen en el objeto (normalmente el centro de masa) y contar el n\'{u}mero de part\'{i}culas (masa = p\'{i}xeles) que componen el objeto en un radio del origen. Para un objeto euclidiano bidimensional, la relaci\'{o}n masa-radio es:
\begin{equation}
M(r) \propto r^{2}
\end{equation}

Por ejemplo, el \'{a}rea de un plano bajo discos de tamaño creciente es proporcional al cuadrado del radio del disco de medici\'{o}n. El exponente es por tanto la dimensi\'{o}n (de hecho, un cuadrado es bidimensional), pero la masa de un objeto fractal incrustado en dos dimensiones cambia con un exponente fraccionario:

\begin{equation}
M(r) \propto r^{D} 
\end{equation}

Y la dimensi\'{o}n fractal $D_\textit{{masa-radio}}$ se obtiene de:

\begin{equation}
D_\textit{masa radio} = \frac{log (M(r))}{log (r)}
\end{equation}

Se calcula como la pendiente de la regresi\'{o}n lineal de $log(M(r))$ de $log(r)$ \cite{Mustafa1996}.


\subsection{M\'{e}todo de masa}

Es una medida similar al m\'{e}todo de conteo de cajas, la diferencia principal es el uso de c\'{i}rculos en lugar de celdas o p\'{i}xeles. Consiste en selecionar un punto de origen en el objeto de estudio (que usualmente es el centro de masas) y contar el n\'{u}mero de part\'{i}culas (masa = c\'{i}rculos) que componen 	el objeto en un radio r desde el origen. Para dos dimesiones la relaci\'{o}n entre masa y radio es:

\begin{equation}
\mu(d) = Ad^D
\end{equation}

Donde $\mu(d)$ es el n\'{u}mero de pixeles en una caja de tamaño d, d es el di\'{a}metro longitud de la caja, A es un variable y D es la dimensi\'{o}n fractal\cite{Mustafa1996}. 

\subsection{Radio de giro}

Existen estructuras que crecen mediante la adici\'{o}n sucesiva de part\'{i}culas al objeto. Esta situaci\'{o}n es com\'{u}n, por ejemplo, en simulaciones de Monte Carlo, donde se registra el n\'{u}mero total de part\'{i}culas dentro de un c\'{u}mulo durante el crecimiento. Por lo tanto, es necesario calcular una cantidad llamada radio de giro  \(R_g(N)\) usando la siguiente expresi\'{o}n:

\begin{equation}
R_g = \sqrt{\frac{1}{N} \sum_{i=1}^{N} \| r_i - r_c \|^2}
\end{equation}

donde \(r_i\) es la distancia de la \(i\)\'{e}sima part\'{i}cula desde el centro de masa del c\'{u}mulo y \(N\) es el n\'{u}mero total de part\'{i}culas en el c\'{u}mulo en la etapa dada del proceso de crecimiento. Luego, se asume que:

\begin{equation}
R_g(N) \approx N^{1/D}
\end{equation}

Por lo tanto, \(1/D\) se puede obtener de la pendiente del gr\'{a}fico de \(\ln(R_g)\) como funci\'{o}n de \(\ln(N)\). Esta relaci\'{o}n corresponde a las siguientes suposiciones: i) en el r\'{e}gimen asint\'{o}tico, \(R_g\) es proporcional linealmente al radio total del c\'{u}mulo, ii) las correcciones debido a efectos de borde se pueden ignorar y iii) la estructura no es un multifractal geom\'{e}trico \cite{Mroczka2012, Vicsek1992}.

\subsection{Multifractalidad}

Las medidas multifactales o fractales est\'{a}n relacionadas con el estudio de una distribuci\'{o}n de cantidades f\'{i}sicas o de otro tipo en un soporte geom\'{e}trico. El soporte puede ser un plano ordinario, la superficie de un volumen, o podr\'{i}a ser en s\'{i} mismo un fractal. En general, un multifractal posee un n\'{u}mero infinito de singularidades de infinitos tipos. La multifractalidad expresa el hecho de que los puntos correspondientes a un tipo determinado de singularidad suelen formar un subconjunto fractal cuya dimensi\'{o}n depende del tipo de singularidad. La idea de que un multifractal puede ser representado en t\'{e}rminos de subconjuntos fractales entrelazados que tienen diferentes exponentes de escala.

\subsection{Geometr\'{i}a fractal en prote\'{i}nas}

Ahora que hemos definido D y explicado algunos m\'{e}todos para obtenerlo, surge naturalmente la pregunta de si el conocimiento de la magnitud de la dimensi\'{o}n tiene alguna utilidad. Las respuestas provienen de varios experimentos en los que se ha empleado para cuantificar algunos aspectos de la morfolog\'{i}a proteica. Dado que la estructura de las prote\'{i}nas tiene una forma tan compleja que solo se puede analizar adecuadamente mediante el enfoque de la geometr\'{i}a fractal \cite{Mustafa1996}. 

Una prote\'{i}na consiste en una cadena de polip\'{e}ptidos compuesta por residuos de amino\'{a}cidos unidos entre s\'{i} por enlaces pept\'{i}dicos. La cadena polipept\'{i}dica o columna vertebral forma un pol\'{i}mero lineal compuesto por unidades repetitivas que son id\'{e}nticas, excepto por los terminales de la cadena \cite{Mustafa1996}.

La estructura primaria de un proteina se conoce como la secuencia de \'{a}cidos am\'{i}nicos y determina la conformaci\'{o}n nativa. Los elementos de la estructura secundaria son las l\'{a}minas a-helices y B-plisadas, mientras que la estructura terciaria se relaciona con las disposiciones espaciales generales de los residuos de amino\'{a}cidos en las prote\'{i}nas. Por lo tanto, la dimensi\'{o}n fractal se puede utilizar para caracterizar las estructuras terciarias de prote\'{i}nas y enzimas. Tambi\'{e}n se pueden describir mediante diferentes modelos fractales y diferentes tipos de dimensiones fractales. Muchos trabajos pioneros han demostrado una y otra vez que la prote\'{i}na que tiene autosimilitud estad\'{i}stica y puede ser manejada por el enfoque fractal. Aunque se pueden construir fractales iterados que son perfectamente autosimilares, la autosimilitud de una cadena de prote\'{i}nas es cierta en lo com\'{u}n solo en un sentido estad\'{i}stico, es decir, la parte no siempre se ver\'{a} exactamente como el todo. Otra diferencia entre las macromol\'{e}culas biol\'{o}gicas y los objetos ideales es que una fractal no es autosimilar en todas las escalas de longitud. Hay l\'{i}mites de tamaño superior e inferior m\'{a}s all\'{a} de los cuales una macromol\'{e}cula ya no es fractal. Por lo que la investigaci\'{o}n fractal sobre enzimas y prote\'{i}nas es actualmente un campo activo \cite{Mustafa1996}. 



\subsection{Medidas basadas en distancia de la similitud estructural de prote\'{i}nas: RMSD}

A pesar de su aparente simplicidad, el problema de cuantificar las diferencias entre dos estructuras de la misma proteína o complejo no es trivial y continúa evolucionando \cite{Kufareva2012}. La Desviaci\'{o}n Cuadr\'{a}tica Media (RMSD) es la medida cuantitativa m\'{a}s com\'{u}nmente usada para la similitud entre dos conjuntos de coordenadas at\'{o}micas superpuestas. Los valores de RMSD se presentan en Ångströms y se calculan mediante:

\begin{equation}
RMSD = \sqrt{\frac{1}{n} \sum_{i=1}^{n} d_i^2}
\end{equation}

donde el promedio se realiza sobre los \(n\) pares de \'{a}tomos equivalentes y \(d_i\) es la distancia entre los dos \'{a}tomos del par \(i\). 

El RMSD se puede calcular para cualquier tipo y subconjunto de \'{a}tomos; por ejemplo, para los \'{a}tomos \(C_{\alpha}\) de toda la prote\'{i}na, para los \'{a}tomos \(C_{\alpha}\) de todos los residuos en un subconjunto espec\'{i}fico (como las h\'{e}lices transmembrana, el sitio de uni\'{o}n o un bucle), todos los \'{a}tomos pesados de un subconjunto espec\'{i}fico de residuos, o todos los \'{a}tomos pesados en un ligando de mol\'{e}cula pequeña. La principal desventaja del RMSD es que est\'{a} dominado por las amplitudes de errores. Dos estructuras que son id\'{e}nticas excepto por la posici\'{o}n de un solo bucle o un extremo flexible generalmente tienen un RMSD global del esqueleto grande y no se pueden superponer eficazmente mediante ning\'{u}n algoritmo que optimice el RMSD global.Una extensi\'{o}n importante de la medida RMSD, es el RMSD ponderado (wRMSD), permite centrarse en subconjuntos at\'{o}micos seleccionados, por ejemplo, minimizando la importancia de las regiones conocidas por ser inherentemente no estructuradas \cite{Kufareva2012}.

\begin{equation}
RMSD = \sqrt{\frac{\sum_{i=1}^{n} w_{i}d_{i}^{2}}{\sum_{i=1}^{n} w_{i}}}
\end{equation}

\clearpage

\subsection{Enfermedades gen\'{e}ticas}

Las tubulinas son las prote\'{i}nas principales que componen los microt\'{u}bulos (estructuras cil\'{i}ndricas huecas que forman parte del citoesqueleto y desempeñan un papel crucial en la divisi\'{o}n celular, el transporte intracelular de ves\'{i}culas y org\'{a}nulos, y la estabilidad celular).
Las tubulinopat\'{i}as son un grupo de enfermedades gen\'{e}ticas raras que afectan la funci\'{o}n de las tubulinas. Las tubulinopat\'{i}as pueden ser causadas por mutaciones gen\'{e}ticas que afectan a los genes que codifican las tubulinas o las prote\'{i}nas asociadas con los microt\'{u}bulos. Estas mutaciones pueden interferir con la funci\'{o}n normal de las tubulinas y los microt\'{u}bulos, lo que resulta en una amplia variedad de manifestaciones cl\'{i}nicas. Las caracter\'{i}sticas cl\'{i}nicas de las tubulinopat\'{i}as pueden incluir anomal\'{i}as del desarrollo del sistema nervioso central, como malformaciones cerebrales, trastornos del desarrollo cognitivo y del habla, retraso en el desarrollo motor, epilepsia y trastornos del movimiento. Tambi\'{e}n pueden afectar otros sistemas del cuerpo, como el sistema ocular y el sistema hematol\'{o}gico. Algunos ejemplos de tubulinopat\'{i}as incluyen el s\'{i}ndrome de lissencefalia tipo 2 (tambi\'{e}n conocido como s\'{i}ndrome de Miller-Dieker), el s\'{i}ndrome de pachygyria liso, y la displasia cortical focal. Dado que las tubulinopat\'{i}as son enfermedades raras, la investigaci\'{o}n en este campo est\'{a} en curso para comprender mejor su patog\'{e}nesis, sus manifestaciones cl\'{i}nicas y para desarrollar enfoques terap\'{e}uticos dirigidos.

\section{Objetivos}
\subsection{Objetivo general}

Determinar si las prote\'{i}nas presentan caracter\'{i}sticas multifractales. De ser as\'{i}, analizar y caracterizar
dicha multifractalidad con el fin de utilizar una o m\'{a}s variantes de la
dimensi\'{o}n fractal como medida(s) caracterizadora(s)
de la estructura y patrones de agregaci\'{o}n de las prote\'{i}nas.


\subsection{Objetivos particulares}

\begin{itemize}

\item Seleccionar un conjunto de prote\'{i}nas de estudio que tengan potencial para aplicaciones posteriores.
Por ejemplo prote\'{i}nas nativas/mutadas, prote\'{i}nas que han mantenido su funci\'{o}n biol\'{o}gica pero que se
han diferenciado estructuralmente, entre otros casos. Se priorizar\'{a}n casos en que se cuente con informaci\'{o}n
experimental (PDB).

\item Adaptar la determinaci\'{o}n de la dimensi\'{o}n fractal de masa a las prote\'{i}nas, a partir de otros casos como geles o materiales ferromagn\'{e}ticos.

\item Analizar la dimension fractal de masa de un conjunto de prote\'{i}nas, partiendo de la definici\'{o}n de multifractalidad y usando los datos experimentales recabados.

\item Analizar si la multifractalidad de una prote\'{i}na se observa con otras formas de determinar la
dimensi\'{o}n fractal de masa, por ejemplo el radio de giro.

\item Explorar el uso potencial de la dimensi\'{o}n fractal de masa como una herramienta para identificar
o resolver entre prote\'{i}nas gen\'{e}ticamente similares pero con diferencias estructurales notables.

\end{itemize}

\clearpage

\section{Metodolog\'{i}a}

\begin{enumerate}
	
\item Obtenci\'{o}n de la estructura tridimensional de la prote\'{i}na: Obtener la estructura tridimensional de la prote\'{i}nas que fue determinda mediante t\'{e}cnicas experimentales como la cristalograf\'{i}a de rayos X o la espectroscopia de resonancia magn\'{e}tica nuclear.

\item C\'{a}lculo de la dimensi\'{o}n fractal de masa: El c\'{a}lculo de la dimensi\'{o}n fractal de masa implica medir c\'{o}mo var\'{i}a el volumen del objeto (en este caso, la prote\'{i}na) con respecto a la escala a la que se est\'{a} observando. Esto se puede hacer utilizando algoritmos espec\'{i}ficos para el c\'{a}lculo de la dimensi\'{o}n fractal de masa como el algoritmo box-counting.

\item An\'{a}lisis de la complejidad geom\'{e}trica de la prote\'{i}na: Una vez calculada la dimensi\'{o}n fractal de masa para la prote\'{i}na, se analizar\'{a} c\'{o}mo var\'{i}a esta dimensi\'{o}n en diferentes escalas. Esto nos permitir\'{i}a caracterizar la complejidad geom\'{e}trica de la prote\'{i}na a diferentes niveles de detalle.

\end{enumerate}

\clearpage

\section{Cronograma}
\begin{table}[hbp!]
\centering
\footnotesize
\setlength{\tabcolsep}{2.0pt}
\begin{tabular}{||p{0.4\linewidth}|c|c|c|c|c|c|c|c|c|c|c|c||}
\hline
\textbf{Actividad} & \multicolumn{5}{c|}{2023} & \multicolumn{7}{c||}{2024}\\
\hline
& Ago & Sep & Oct & Nov & Dic & Ene & Feb & Mar & Abr & May & Jun & Jul\\
\hline
Materias del primer semestre & X & X & X & X & X & & & & & & &  \\
\hline
LII Winter Meeting on Statistical Physics & & & & & & X & & & & & & \\
\hline
Curso de Qu\'{i}mica Cu\'{a}ntica en \'{a}tomos & & & & & & X & X & X & X & X & X & \\
\hline
Materias del segundo semestre &  &  &  &  &  & X & X & X & X & X & X &\\
\hline
Revisi\'{o}n bibliogr\'{a}fica & X & X & X & X & X & X & X & X & X & X & X & X \\
\hline
Presentaci\'{o}n del protocolo &  &  &  &  &  &  &  &  &  &  & X & \\
\hline
Preparaci\'{o}n de sistemas de estudio  &  &  &  &  &  &  &  &  &  &  &  & X \\
\hline
& \multicolumn{5}{c|}{2024} & \multicolumn{7}{c||}{2025}\\\hline
& Ago & Sep & Oct & Nov & Dic & Ene & Feb & Mar & Abr & May & Jun & Jul\\
\hline
XXII RMFQT &  &  &  & X &  &  &  &  &  &  &  &\\
\hline
XXlX Simposio interno FCQ - ICUAP &  &  &  & X &  &  &  &  &  &  &  &\\
\hline
LIII Winter Meeting on Statistical Physics &  &  &  &  &  & X &  &  &  &  &  &\\
\hline
Revisi\'{o}n bibliogr\'{a}fica & X & X & X & X & X & X & X & X & X & X & X &\\
\hline
Redacci\'{o}n de tesis &  &  &  &  &  & X & X & X & X &  &  &\\
\hline
Revisi\'{o}n de tesis &  &  &  &  &  &  &  &  & X & X & X & \\
\hline
Examen &  &  &  &  &  &  &  &  &  &  &  & X\\
\hline
%Actividad &  &  &  &  &  &  &  &  &  &  &  &  \\
\hline
\end{tabular}
\end{table}


\clearpage

\bibliographystyle{plain}
\bibliography{bibliografia}

\end{document}

